%----------------------------------------------------------------------
% 附錄頁
%----------------------------------------------------------------------

\begin{Appx}{Appendix}

\section{Wave equations from Maxwell's equations in free space}
 \label{ch:Maxwell} 
In the electrodynamics framework, Maxwell’s equations can be written as
\begin{align} 
\boldsymbol{\nabla}\cdot\boldsymbol{E}(\boldsymbol{r}, t)&=\frac{\rho}{\epsilon_0}= 0, \label{max1}\\ 
\boldsymbol{\nabla}\cdot\boldsymbol{B}(\boldsymbol{r}, t)&= 0,\label{max2} \\ 
\boldsymbol{\nabla}\times\boldsymbol{E}(\boldsymbol{r}, t) &=  -\frac{\partial }{\partial t}\boldsymbol{B}(\boldsymbol{r}, t),\label{max3}\\
\boldsymbol{\nabla}\times\boldsymbol{B}(\boldsymbol{r}, t) &=\mu_0\mathbf{J} + \frac{1}{c^2}\frac{\partial }{\partial t}\boldsymbol{E}(\boldsymbol{r}, t)= \frac{1}{c^2}\frac{\partial }{\partial t}\boldsymbol{E}(\boldsymbol{r}, t).\label{max4}
\end{align}
Here $\boldsymbol{\nabla}\equiv\frac{\partial}{\partial x}\hat{\boldsymbol{x}}+\frac{\partial}{\partial y}\hat{\boldsymbol{y}}+\frac{\partial}{\partial z}\hat{\boldsymbol{z}}$ represents the gradient operator along the three spatial coordinates. 
These equations can also be expressed in terms of electromagnetic potentials, which will be useful later when discussing light–matter interaction. First, the divergence-free condition of Eq.~(\ref{max2}) allows the magnetic field to be written in terms of a \textbf{vector potential} $\boldsymbol{A}(\boldsymbol{r}, t)$ as
\begin{equation} 
\boldsymbol{B}(\boldsymbol{r}, t)=\boldsymbol{\nabla}\times\boldsymbol{A}(\boldsymbol{r}, t). \label{BcrulA}
\end{equation}
Substituting this expression into Faraday's law (Eq.~(\ref{max3})) yields a quantity whose curl vanishes. This allows the introduction of a \textbf{scalar potential} $\phi(\boldsymbol{r}, t)$ such that
\begin{equation} 
\boldsymbol{\nabla}\phi(\boldsymbol{r}, t)=\left(\boldsymbol{E}(\boldsymbol{r}, t)+\frac{\partial }{\partial t}\boldsymbol{A}(\boldsymbol{r}, t)\right).
\end{equation}
Therefore, the electric field can be expressed in terms of the scalar and vector potentials as
\begin{equation} 
\boldsymbol{E}(\boldsymbol{r}, t)=-\boldsymbol{\nabla}\phi(\boldsymbol{r}, t)-\frac{\partial }{\partial t}\boldsymbol{A}(\boldsymbol{r}, t).
\label{gaugeF}
\end{equation}
This expression automatically satisfies both Eq.~(\ref{max2}) and Faraday's law (Eq.~(\ref{max3})). Finally, Gauss's law (Eq.~(\ref{max1})) and Ampère's law with Maxwell's correction (Eq.~(\ref{max4})) can be rewritten as two coupled wave equations for the electromagnetic potentials:
\begin{align} 
&\nabla^2\phi(\boldsymbol{r}, t)+\frac{\partial }{\partial t}\left(\boldsymbol{\nabla}\cdot\boldsymbol{A}(\boldsymbol{r}, t)\right)=0\\
&\nabla^2\boldsymbol{A}(\boldsymbol{r}, t)-\frac{1}{c^2}\frac{\partial^2 }{\partial t^2}\boldsymbol{A}(\boldsymbol{r}, t)-\boldsymbol{\nabla}\left(\boldsymbol{\nabla}\cdot\boldsymbol{A}(\boldsymbol{r}, t)+\frac{1}{c^2}\frac{\partial }{\partial t}\phi(\boldsymbol{r}, t)\right)=0. 
\end{align}



\section{scalar Helmholtz equation}
 \label{ch:SHE} 

We begin with the Lorenz gauge. First, we assume a \textbf{monochromatic wave} ansatz,$\boldsymbol{A}^L(\boldsymbol{r}, t)=\text{Re}\left[\boldsymbol{A}^L(\boldsymbol{r}) e^{-i\omega t}\right]$,which transforms Eq.~(\ref{Lwave2}) into the vector Helmholtz equation
\begin{equation} 
\nabla^2\boldsymbol{A}^L(\boldsymbol{r})+q_0^2\boldsymbol{A}^L(\boldsymbol{r})=0. 
\label{VH}
\end{equation}

Next, we introduce the \textbf{transverse-wave} ansatz,$\boldsymbol{A}^L(\boldsymbol{r})=\hat{\boldsymbol{\varepsilon}}A^L(\boldsymbol{r})$, where $\hat{\boldsymbol{\varepsilon}}$ is a polarization unit vector. Substituting this form into Eq.~(\ref{VH}) reduces the equation to the scalar Helmholtz equation
\begin{equation} 
\nabla^2A^L(\boldsymbol{r})+q_0^2A^L(\boldsymbol{r})=0.
\label{SH}
\end{equation}

\end{Appx}